\subsection{An allowance affects capital and exposure in different ways}
\label{sec:bank-prudential-connection}

The illustrative ratio above shows where a prudential calculation belongs.
Applying a particular rule requires a further step: reconcile its numerator
and denominator to the same bank's accounting records. An allowance already
reduces accounting equity through expense, yet a specified portion may enter
Tier~2 capital. The regulatory loan exposure can meanwhile remain gross of
that allowance. Using the net loan balance for every purpose would erase
these distinctions.

Consider a separate synthetic US national bank, with amounts in millions of
US dollars. Its conditional reporting window is January--March 2026. For
this calculation, freeze the retained 1 January 2026 national-bank rules;
this assumption does not establish the rules actually applicable to a
31 March filing. The bank has no subsidiaries, minority interests, AOCI,
deferred taxes, other capital instruments or deductions apart from goodwill.
It is outside the advanced-approaches, Category~III, market-risk and
community-bank leverage-ratio regimes, and has no transition or special
supervisory adjustment. Its common stock is supplied as satisfying all
thirteen conditions in section 3.20(b)(1). These are explicit case facts,
not conclusions drawn from a security's name
\citep[12~CFR~3.20(b)(1)--(2)]{occcapitaleligibility2026}.

Settlement cash is a direct Federal Reserve balance. The loans are ordinary,
unsecured, performing corporate exposures measured at amortized cost: no
purchased credit-deteriorated loans, nonaccrual, ninety-day delinquency,
HVCRE, securities, derivatives or credit mitigation. The bank has adopted
CECL. Its allowance amounts are supplied accounting measurements, not
estimates made here, and all qualify as adjusted allowances for credit
losses, or AACL. There is no allocated transfer risk reserve or off-balance
credit exposure. These restrictions are consequential: a deposit at another
commercial bank and a nonaccrual corporate loan need different exposure
analysis \citep[12~CFR~3.2, definitions of AACL and carrying value]{occcapitaldefinitions2026}.

The opening book contains cash 30, gross loans 100, goodwill 2 and allowance
0.8 against deposits 110, common stock 20 and retained earnings 1.2.
Assets net of allowance are \(30+100+2-0.8=131.2\), so equity is 21.2.
Originate a loan of 20 by crediting the borrower's deposit, transfer 15
externally, receive earned loan interest 6, and credit deposit interest 2.
Remeasure the supplied allowance to 2.4, recording credit-loss expense 1.6.
The income statement reports \(6-2-1.6=2.4\), which a separate closing entry
transfers to retained earnings. Cash becomes 21, loans 120, goodwill 2,
allowance 2.4 and deposits 117; net assets 140.6 equal deposits 117 plus
equity 23.6. The cash bridge is \(30-15+6=21\). The loan origination,
deposit-interest credit and allowance change do not move settlement cash.

\subsubsection*{Construct the exposure base before applying the allowance cap}

For these ordinary loans, the regulatory carrying-value definition does not
subtract the associated allowance. The exposure amount is therefore the
gross loan balance. Section 3.31 multiplies that amount by the applicable
weight; section 3.32 assigns 100\% to the stated corporate exposures and
zero to the direct central-bank balance. Goodwill is deducted from CET1
and excluded from risk-weighted assets. With no associated deferred tax
liability, its deduction is the full 2
\citep[12~CFR~3.31(a)--(b)]{occrwamechanics2026}
\citep[12~CFR~3.32(a)(1)(i)(A), (f)(1), (k)]{occriskweights2026}
\citep[12~CFR~3.22(a)(1), (e), (g)]{occcapitaldeductions2026}.

Let \(G\) be gross loans, \(A\) eligible AACL and \(B\) the risk-weighted
base used to limit the allowance admitted to Tier~2. For this restricted
book, \(B=G\). The capital rule caps that admission at 1.25\% of the
relevant base. The official reporting instructions make the order concrete:
Part~I item~42 uses Part~II item~26 for the cap; Part~II subsequently
subtracts the excess allowance in arriving at its total. Define admitted
allowance \(T\), excess \(X\) and final standardized risk-weighted assets
\(D\) accordingly:
\[
 T=\min\{A,0.0125B\},\qquad X=A-T,\qquad D=B-X.
\]
The last equality uses this example's exclusions. In a broader book, the
cap base and pre-excess standardized denominator need not coincide. The
instructions exclude section 22(a) deductions such as goodwill from the
cap base but include specified section 22(c)--(f) deductions, which remain
excluded from the final standardized denominator. The case has none of
those other deductions or market-risk amounts
\citep[12~CFR~3.20(d)(3)]{occcapitaleligibility2026}
\citep[Part~I, item~42, pp.~RC-R-27--28; Part~II, items~26 and 28--31,
pp.~RC-R-109--110]{ffiec051capital2024,ffiec051rwa2024}.

These retained instructions are the December 2024 files, with March 2024
revision footers on the cap pages. They corroborate the ordering in the
specified historical rule calculation; they do not supply a complete
amendment history for an actual later filing. The source date and the
conditional reporting date remain separate.

Initially \(B=100\), the cap is 1.25 and all the allowance 0.8 enters
Tier~2, leaving \(D=100\). At the final pre-runoff state, \(B=120\), the
cap is 1.5, and excess allowance is \(2.4-1.5=0.9\). Thus
\(D=120-0.9=119.1\). CET1 is equity less goodwill, \(23.6-2=21.6\).
There is no additional Tier~1 capital, so Tier~1 is also 21.6; total
capital is \(21.6+1.5=23.1\). The following reconciliation retains the
different measures rather than giving them one label.
\begin{center}
\begin{tabular}{lrr}
\toprule
Measure & Opening & Before runoff\\
\midrule
Accounting equity & 21.2 & 23.6\\
Goodwill deduction & 2 & 2\\
CET1 and Tier 1 & 19.2 & 21.6\\
Eligible allowance in Tier 2 & 0.8 & 1.5\\
Total capital & 20 & 23.1\\
Allowance-cap base & 100 & 120\\
Excess allowance & 0 & 0.9\\
Standardized risk-weighted assets & 100 & 119.1\\
\bottomrule
\end{tabular}
\end{center}
CET1 and Tier~1 ratios are \(21.6/119.1=72/397\), about 18.14\%;
the total-capital ratio is \(23.1/119.1=77/397\), about 19.40\%.
These are conditional ratios, not a determination of payout permission or
all applicable capital requirements
\citep[12~CFR~3.10(b)(1)--(3)]{occcapitalratios2026}.

Several plausible substitutions give the wrong quantity. Net loans 117.6
are the accounting amount, not this regulatory exposure. Deducting the
full allowance from risk-weighted assets loses the rule's distinction
between admitted and excess allowance. Deducting it again from CET1
counts an expense twice. Including goodwill in the exposure base reverses
its required exclusion. Finally, using \(D\) to calculate its own
allowance cap gives the circular equation \(T=(117.6+T)/80\), hence
\(T=117.6/79\), rather than the required 1.5. Solving that equation
accurately cannot make it the specified calculation.

\subsubsection*{A capital ratio does not settle the next cash payment}

Now suppose 25 of deposits runs off. The supplied funding policy draws 4
within an available facility limit of 10, then pays 25. Cash becomes zero,
deposits 92 and borrowing 4. Net assets 119.6 still equal liabilities 96
plus equity 23.6. Capital and risk-weighted assets are unchanged: the
disappearing central-bank cash has zero weight, and no additional income
or expense is assumed. The facility is borrowing by this bank, not an
off-balance extension of credit by it to a customer.

If instead 40 is demanded, the same facility supplies at most 10. The bank
pays 31 and leaves 9 unpaid, retained within its deposit liability of 86.
Borrowing is 10 and cash zero. The capital ratios remain the same, although
the payment obligation has failed. No capital or cash adjustment may be
invented to remove that default. The two runoff outcomes are alternative
histories of this bank, not balances to be combined.

The ordinary leverage ratio asks for another denominator: average total
consolidated assets as reported for the relevant window, less specified
Tier~1 deductions. Neither closing assets 140.6 before runoff nor 119.6
after it determines that average. If a separate, properly matched input
supplies average assets 145, the denominator is \(145-2=143\) and the
ratio is \(21.6/143=108/715\), about 15.10\%. Without that input, the
leverage ratio remains unresolved while the supported risk-based ratios
remain available \citep[12~CFR~3.10(b)(4)]{occcapitalratios2026}.

The proposed \texttt{RegulatoryAssessment} therefore references the exact
bank, book, perimeter, accounting state, information view, rule snapshot,
capital eligibility, exposure classifications, deductions, allowance
admission and reporting window. Average-assets evidence is a separate
dependency. Its result names each numerator and denominator and records
which outputs remain unresolved. Unknown stock eligibility blocks the
dependent capital amounts; an incorrect perimeter rejects the input; a
missing average blocks leverage alone. The assessment creates no accounting
entry. A changed allowance, instrument, exposure class, average or rule
requires reassessing the affected outputs while retaining the historical
result under its original assumptions.

